\section{Marco Teórico}
\label{sec:marco}

\subsection{Heladas}
[Define helada y sus tipos (radiativas y advectivas), y cómo afectan a los cultivos altoandinos. Cita tus fuentes, por ejemplo SENAMHI.]

\subsection{Series temporales}
[Define serie temporal, estacionalidad, tendencia y la forma de construir variables rezagadas para predecir a 24 y 48 horas.]

\subsection{Algoritmos de aprendizaje automático}
[Regresión logística; Random Forest \cite{breiman2001}; XGBoost \cite{chen2016}; LSTM \cite{hochreiter1997}. Explica brevemente cada uno.]

\subsection{Métricas para clases desbalanceadas}
Dado que la helada es un evento poco frecuente, la exactitud global no es suficiente y se emplean las métricas de la Tabla~\ref{tab:metricas}.

\begin{table}[ht]
\centering
\caption{Métricas de evaluación para la clase helada}
\label{tab:metricas}
\begin{tabular}{@{}p{3cm}p{11cm}@{}}
\toprule
\textbf{Métrica} & \textbf{Qué mide} \\
\midrule
Recall & Proporción de heladas reales que el modelo logra detectar. \\
Precisión & Proporción de alertas emitidas que corresponden a heladas reales. \\
F1 & Media armónica entre precisión y recall. \\
AUC-PR & Área bajo la curva precisión--recall; resume el desempeño para distintos umbrales. \\
MCC & Coeficiente de correlación de Matthews; equilibra los cuatro valores de la matriz de confusión. \\
\bottomrule
\end{tabular}\\
\textit{Nota.} Elaboración propia.
\end{table}